\documentclass[UTF8,a4paper,12pt]{ctexart}

\usepackage[left=2.5cm,right=2.5cm,top=2.6cm,bottom=2.6cm]{geometry}
\usepackage{amsmath,amssymb,bm}
\usepackage{graphicx}
\usepackage{booktabs}
\usepackage{array}
\usepackage{float}
\usepackage{tabularx}
\usepackage{listings}
\usepackage{xcolor}
\usepackage{hyperref}
\usepackage{enumitem}
\usepackage{url}
\usepackage{makecell}

\graphicspath{{fig/}{report/fig/}}

\newcommand{\safeincludegraphics}[2][]{%
  \IfFileExists{#2}{%
    \includegraphics[#1]{#2}%
  }{%
    \IfFileExists{fig/#2}{%
      \includegraphics[#1]{fig/#2}%
    }{%
      \IfFileExists{report/fig/#2}{%
        \includegraphics[#1]{report/fig/#2}%
      }{%
        \fbox{\parbox{0.82\textwidth}{\centering Missing figure: \texttt{\detokenize{#2}}}}%
      }%
    }%
  }%
}

\hypersetup{colorlinks=true, linkcolor=black, citecolor=black, urlcolor=blue}
\setlist[itemize]{leftmargin=1.8em}
\setlist[enumerate]{leftmargin=1.8em}
\renewcommand{\arraystretch}{1.15}

\lstset{
  basicstyle=\small\ttfamily,
  keywordstyle=\color{blue!70!black},
  commentstyle=\color{green!35!black},
  stringstyle=\color{orange!60!black},
  frame=single,
  numbers=left,
  numberstyle=\tiny,
  breaklines=true,
  showstringspaces=false,
  columns=fullflexible
}

\title{PageRank Report with CSR and Block Matrix Optimization}
\author{bigdatacompute}
\date{\today}

\begin{document}
\maketitle
\tableofcontents
\newpage

\section{第 1 章 任务背景与数据集}


\subsection{任务背景}


本次作业要求在给定有向图数据集 \texttt{Data.txt} 上计算 PageRank 分数，并输出排名前 100 的节点及其分数到 \texttt{Res.txt}。输入文件每一行包含两个整数 \texttt{FromNodeID ToNodeID}，表示一条从源节点指向目标节点的有向边。作业明确要求程序需要迭代至收敛，必须处理 dead-end 与 spider-trap，并且至少报告阻尼系数 \texttt{beta = 0.85} 时的 Top-100 结果。


本题同时考察 PageRank 正确性与内存优化。实现采用 CSR 压缩图结构，并结合按目标节点分块与 \texttt{memmap} 控制迭代阶段的活动工作集，以满足 \texttt{80 MB} 内存与 \texttt{60 s} 时间约束。


\subsection{数据集规模}


使用 \texttt{analyze\_dataset.py} 对 \texttt{Data.txt} 做一次完整扫描，得到如下统计结果：


\begin{center}
\small
\begin{tabularx}{\textwidth}{>{\raggedright\arraybackslash}X >{\raggedright\arraybackslash}X}
\toprule
指标 & 数值 \\
\midrule
节点数 \texttt{N} & \texttt{10000} \\
边数 \texttt{M} & \texttt{150000} \\
平均出度 & \texttt{15.00} \\
平均入度 & \texttt{15.00} \\
最大出度 & \texttt{36} \\
最大入度 & \texttt{32} \\
最大总度数 & \texttt{57} \\
dead-end 节点数 & \texttt{1500} \\
dead-end 占比 & \texttt{15.00\%} \\
弱连通分量数 & \texttt{501} \\
最大弱连通分量节点数 & \texttt{9500} \\
最大弱连通分量占比 & \texttt{95.00\%} \\
孤立节点数 & \texttt{500} \\
自环数 & \texttt{15} \\
重复边数 & \texttt{0} \\
\bottomrule
\end{tabularx}
\end{center}


从规模上看，本数据集节点数为一万，边数为十五万，图整体是稀疏图。若使用稠密邻接矩阵，单是 \texttt{10000} $\times$ \texttt{10000} 的 \texttt{float32} 矩阵就需要约 \texttt{381 MB} 以上，不可能满足内存限制；而 CSR 只需要保存真实存在的边和行指针，核心数组规模可以压到 MB 级以内。这也是选择 CSR 作为主存储结构的直接原因。


\subsection{度分布与 dead-end 特征}


度分布如图~\ref{fig:degree-distribution} 所示。平均入度和平均出度都等于 \texttt{15}，这是由总边数与节点数的比例决定的。但最大出度为 \texttt{36}、最大入度为 \texttt{32}，说明节点连接并不完全均匀，少数节点拥有更高的链接集中度。


\begin{figure}[H]
\centering
\safeincludegraphics[width=0.86\textwidth]{degree_distribution.png}
\caption{数据集入度与出度分布。}
\label{fig:degree-distribution}
\end{figure}



数据集中还有 \texttt{1500} 个 dead-end 节点，占 \texttt{15.00\%}。这些节点没有出边，如果在迭代中直接忽略，概率质量会逐渐流失，PageRank 向量不再归一化。实现中采用统一的质量补偿公式：每轮把 dangling mass 按 \texttt{beta * dangling\_mass / N} 均分回所有节点，再叠加随机跳转项 \texttt{(1 - beta) / N}。图~\ref{fig:deadend-strategy} 对比了不同策略的效果。


\begin{figure}[H]
\centering
\safeincludegraphics[width=0.86\textwidth]{deadend_strategy_comparison.png}
\caption{不同 dead-end 处理策略对概率质量和 PageRank 结果的影响。}
\label{fig:deadend-strategy}
\end{figure}



\subsection{连通性观察}


按无向图视角统计弱连通分量，数据集中共有 \texttt{501} 个弱连通分量，其中最大弱连通分量包含 \texttt{9500} 个节点，占总节点数 \texttt{95.00\%}。同时存在 \texttt{500} 个孤立节点。这个结构说明图的主体部分高度集中，但也存在一批完全孤立或与主体无连接的小分量。


对 PageRank 而言，弱连通分量较多本身并不会导致算法失效，因为 teleport 机制会让随机游走在全图范围内保持可达性；但它会影响低排名节点的分数分布，也让 dead-end 补偿更重要。Top-10 排名主要受主体连通分量中的高入度或高质量来源节点影响。


\subsection{本章小结}


本数据集是一个典型稀疏有向图：\texttt{N = 10000}、\texttt{M = 150000}、平均度为 \texttt{15}。CSR 能显著降低图存储成本，Block Matrix 则用于进一步控制迭代过程中的内存活动范围。数据集中 \texttt{15.00\%} 的 dead-end 节点和 \texttt{501} 个弱连通分量说明，正确处理 dangling mass 与 teleport 不是附加细节，而是 PageRank 收敛和结果可信的基础。

\section{第 2 章 算法原理}


\subsection{PageRank 的随机游走视角}


PageRank 是经典的网页排序算法，最早由 Page 等人在技术报告中系统提出，并在 Brin 与 Page 的大规模 Web 搜索引擎论文中得到应用；其随机游走和矩阵解释也可参考 Langville 与 Meyer 的专著\cite{page1999pagerank,brin1998anatomy,langville2006google}。PageRank 可以理解为一个“随机浏览者”在有向图上的长期访问概率分布。图中每个节点表示一个页面，边 \texttt{u -> v} 表示页面 \texttt{u} 上存在指向 \texttt{v} 的超链接。若浏览者当前位于节点 \texttt{u}，则有两种行为：


\begin{enumerate}
\item 以概率 \texttt{beta} 沿着 \texttt{u} 的一条出边继续访问；
\item 以概率 \texttt{1 - beta} 发生随机跳转，直接跳到任意一个节点。
\end{enumerate}


如果记第 \texttt{t} 轮时的排名向量为 $r^{(t)}$，图共有 \texttt{N} 个节点，节点 \texttt{u} 的出度为 \texttt{L(u)}，则节点 \texttt{v} 的 PageRank 满足


\[ PR(v)=\frac{1-\beta}{N}+\beta \sum_{u \to v}\frac{PR(u)}{L(u)}. \]


公式由两部分组成：$(1-\beta)/N$ 是随机跳转带来的均匀基础概率，保证任意节点都有被访问的机会；另一部分来自入边源点的贡献，每个源点 \texttt{u} 把自己的权重均分给所有出边终点。


若把图写成转移矩阵 \texttt{M}，其中


\[
M[u,v]=
\begin{cases}
\frac{1}{L(u)}, & u \to v,\\
0, & \text{otherwise}.
\end{cases}
\]


则幂迭代形式可写为


\[ r^{(t+1)}=\beta M^{T}r^{(t)}+\frac{1-\beta}{N}\mathbf{1}. \]


课程实现中采用的就是这一迭代思想：从均匀分布 $r^{(0)} = (1/N, \dots, 1/N)$ 出发，不断重复更新，直到相邻两轮向量变化足够小。


\subsection{dead-end 的成因与影响}


dead-end（悬挂节点）指的是出度为 \texttt{0} 的节点。对这类节点，普通转移矩阵无法定义“把概率往哪里传出去”，因为它没有任何出边。若仍然直接套用基础公式，则这一行会全为 \texttt{0}，结果是到达 dead-end 的概率质量在下一轮迭代中被“吞掉”。随着迭代进行，排名向量的总和会持续下降，最终不再是合法的概率分布。


在实际数据集中，dead-end 并不少见——很多节点只被引用、不再向外链接。如果不处理，概率质量会逐轮泄漏，PageRank 向量总和持续下降，最终不再是合法的概率分布。而且如果团队成员各自用不同策略，实验结果就无法对齐，所以必须统一处理。


本项目冻结接口明确要求全队统一使用“质量补偿法”。设


\[ D=\{i \mid \mathrm{out\_deg}[i]=0\}, \qquad \mathrm{dangling\_mass}^{(t)}=\sum_{i \in D} r_i^{(t)}. \]


则每一轮迭代时，把这部分丢失的概率质量重新平均分配到所有节点，即


\[ r^{(t+1)}[v] = \frac{1-\beta}{N} + \beta \frac{\mathrm{dangling\_mass}^{(t)}}{N} + \beta \sum_{u \to v,\, L(u)>0}\frac{r^{(t)}[u]}{L(u)}. \]


这样处理后，dead-end 节点的权重不会消失，而是像“随机再投放一次”那样回流到全图，向量总质量可以维持在 1 附近。


\subsection{spider-trap 的成因与影响}


spider-trap 指一个内部强连通、但几乎不向外部连边的子图。例如若一组节点之间彼此大量互链，而对外没有出口，那么随机游走一旦进入这部分区域，就会长期被困在其中。若没有随机跳转项，最终几乎全部权重都会被这类结构吸走，得到的排序会高度偏置。


spider-trap 和 dead-end 的区别在于：dead-end 导致概率无法继续传播，spider-trap 则让概率过度集中在局部封闭结构中。两者都会破坏全局随机游走分布，导致 PageRank 结果失真。


PageRank 中的 teleport 项


\[ \frac{1-\beta}{N}\mathbf{1} \]


正是解决 spider-trap 的关键。因为无论当前处于哪个节点，浏览者都有 \texttt{1 - beta} 的概率跳出当前局部结构，回到全图范围重新采样。因此，spider-trap 不需要额外单独设一套修正规则，只要正确实现 teleport，并且 dead-end 也被统一补偿，就能保证收敛到稳定分布。


\subsection{收敛性的理论依据}


加入随机跳转后，实际迭代使用的矩阵（常称为 Google Matrix）在 $0 < \beta < 1$ 时满足以下性质：


\begin{enumerate}
\item 每一列对应的概率质量是可控的；
\item 任意节点之间都存在非零概率互达；
\item 矩阵是 primitive 的，因此存在唯一的主特征向量。
\end{enumerate}


因此无论初始向量怎么选，幂迭代都会收敛到唯一的稳定分布。本实验固定采用


\begin{itemize}
\item \texttt{beta = 0.85}
\item \texttt{eps = 1e-8}
\end{itemize}


\texttt{beta} 太小会稀释图结构信息，太大则对 spider-trap 更敏感、收敛更慢，\texttt{0.85} 是比较常用的折中值。\texttt{eps = 1e-8} 在精度和时间之间取得了平衡，多轮实验中 Top-100 排名保持稳定。


\subsection{收敛判据的选择}


本项目采用


\[ \lVert r^{(t+1)}-r^{(t)} \rVert_1 < \varepsilon \]


作为收敛判据。选用 L1 范数而不是 L2 或最大范数，主要是因为 L1 直接表示全图概率质量的总移动量，物理意义比较清楚，而且对大量小幅变化也比较敏感。此外 L1 判据便于不同实验设置之间进行一致比较。


实现上也比较简单：构造一个差值缓冲区，原地做减法、取绝对值、求和就行，不需要平方或开方。


\subsection{本章小结}


PageRank 求的是随机游走意义下的稳态概率分布，而不只是简单统计入边数量。dead-end 会造成概率质量流失，需要通过质量补偿修正；spider-trap 会造成局部质量滞留，依赖 teleport 项抑制。只要 Google Matrix 构造正确，幂迭代就能收敛到唯一解。第 3 章进一步说明如何将该公式落实为 CSR 稀疏实现。

\section{第 3 章 算法实现细节}


\subsection{总体实现思路}


本项目 A 侧的核心实现围绕三个接口展开：


\begin{enumerate}
\item \texttt{load\_graph(path)}：读取 \texttt{Data.txt}，构造按源点分组的 CSR；
\item \texttt{power\_iteration(...)}：在 CSR 上执行幂迭代；
\item \texttt{dump\_top\_k(path, ranks, k)}：按约定格式输出结果。
\end{enumerate}


实现需要同时满足接口统一、内存受控、实验可复现三项要求，整条链路遵循两个原则：


\begin{itemize}
\item 所有索引数组统一为 \texttt{np.int32}；
\item 所有排名相关浮点数组统一为 \texttt{np.float32}。
\end{itemize}


这样可以保证内存可控，同时让不同模块使用统一的数据约定。


\subsection{节点编号与 ID 映射方案}


课程给定的边文件只有两列：\texttt{FromNodeID ToNodeID}。理论上存在两种可能：


\begin{enumerate}
\item 节点编号本身已经是从 \texttt{0} 开始的连续编号；
\item 节点编号是稀疏的，需要在程序内部做压缩映射。
\end{enumerate}


但是这里有一个实际工程细节：如果文件里某些编号从未出现在任何边中，它们可能代表“完全孤立节点”。若简单按照“见到的节点集合”做紧凑映射，就会把这些节点直接丢掉，导致节点总数 \texttt{N} 变小，进而改变 PageRank 分布。


因此实现采用折中策略：


\begin{enumerate}
\item 先流式读入原始边表，只保存原始 \texttt{src/dst}；
\item 若观察到编号从 \texttt{0} 起步且整体编号密度较高，则保留 \texttt{0..max\_id} 的编号空间，把缺席编号视为孤立点；
\item 若编号并非从 \texttt{0} 起步，或跨度远大于已观测节点数，则退化为紧凑重映射。
\end{enumerate}


这样既能兼顾课程数据集中的孤立点，又能在明显非连续编号场景下完成压缩。由于冻结接口没有把 \texttt{id\_map} 暴露给 \texttt{dump\_top\_k}，实现中额外维护了一个模块级反向映射 \texttt{\_LAST\_NODE\_IDS}，供输出阶段恢复原始 NodeID。


\subsection{CSR 存储布局}


本实现使用经典的 Compressed Sparse Row（CSR）布局，但含义按“源节点的出边列表”组织：


\begin{itemize}
\item \texttt{row\_ptr[u] : row\_ptr[u+1]} 表示源点 \texttt{u} 的出边区间；
\item \texttt{col\_idx[row\_ptr[u] : row\_ptr[u+1]]} 存放这些出边的目标点；
\item \texttt{out\_deg[u] = row\_ptr[u+1] - row\_ptr[u]}。
\end{itemize}


它的优点是结构紧凑、易于序列化，并适合与分块迭代模块对接。相较于稠密矩阵 \texttt{N x N} 的存储方式，CSR 只为真实存在的边分配空间，其空间复杂度为


\[ O(N + M). \]


其中 \texttt{N} 是节点数，\texttt{M} 是边数。对于当前 \texttt{10000} 节点、\texttt{150000} 边的数据集，CSR 版本内存占用显著低于稠密方案。


建图过程分两步：


\begin{enumerate}
\item 统计每个源点的出度，得到 \texttt{out\_deg}；
\item 对 \texttt{out\_deg} 做前缀和，得到 \texttt{row\_ptr}，再借助写指针把目标点写入 \texttt{col\_idx}。
\end{enumerate}


这样不需要对全部边做额外排序，只需一次线性遍历即可完成压缩。


\subsection{\texttt{load\_graph} 的工程实现}


\texttt{load\_graph} 不直接把整个文本读成一个大字符串，而是逐行扫描。这样做有三个好处：


\begin{enumerate}
\item 格式校验可以精确到行号；
\item 不依赖 \texttt{numpy} 的隐式 dtype 推断；
\item 便于后续在读入后立刻释放临时边缓冲。
\end{enumerate}


核心流程如下：


\begin{lstlisting}
load_graph(path):
    逐行读取 (raw_src, raw_dst)
    校验：两列、整数、非负、可落入 int32
    缓存原始边表
    根据编号分布决定：
        保留 0..max_id
        或构造紧凑映射
    统计 out_deg
    前缀和生成 row_ptr
    按写指针填充 col_idx
    删除临时边缓冲并 gc.collect()
    返回 row_ptr, col_idx, out_deg, N
\end{lstlisting}


这里特别注意两点：


\begin{enumerate}
\item 所有返回数组必须已经是 \texttt{np.int32}，不能依赖调用方再做一次转换；
\item 临时边缓冲在建好 CSR 后立即释放，避免高峰内存叠加。
\end{enumerate}


\subsection{幂迭代主循环}


\texttt{power\_iteration} 的输入是 CSR 三元组与参数 \texttt{(beta, eps)}，输出为 \texttt{(ranks, iters, delta)}。实现没有使用任何现成 PageRank API，而是直接按照冻结公式做原地更新。


伪代码如下：


\begin{lstlisting}
power_iteration(row_ptr, col_idx, out_deg, N, beta, eps):
    r      <- [1/N] * N
    r_new  <- empty(N)
    delta_buf <- empty(N)
    live_mask <- out_deg > 0
    dead_mask <- out_deg == 0
    inv_out_deg[live_mask] <- 1 / out_deg[live_mask]

    repeat t = 1 .. max_iter:
        dangling_mass <- sum(r[dead_mask])
        base <- (1 - beta) / N + beta * dangling_mass / N
        r_new.fill(base)

        for each live source node u:
            contrib <- beta * r[u] / out_deg[u]
            for each edge u -> v:
                r_new[v] += contrib

        delta_buf <- abs(r_new - r)
        delta <- sum(delta_buf)
        if delta < eps:
            return r_new, t, delta

        swap(r, r_new)
\end{lstlisting}


实际代码中，为了减少临时对象分配，做了几处工程化处理：


\begin{enumerate}
\item \texttt{r\_new}、\texttt{delta\_buffer} 只申请一次，循环内复用；
\item \texttt{inv\_out\_deg} 预先计算，避免每轮重复除法；
\item \texttt{np.add.at} 用于把单个源点的贡献散射到多个目标点；
\item 差值范数通过缓冲区原地计算，避免每轮新建一份差值数组。
\end{enumerate}


\subsection{dead-end 补偿的工程实现}


dead-end 的统一实现点只有一个：每轮开始先求


\[ \mathrm{dangling\_mass}=\sum_{i:\mathrm{out\_deg}[i]=0} r[i]. \]


然后把


\[ \beta \cdot \mathrm{dangling\_mass} / N \]


直接并入所有节点共享的 \texttt{base} 项。这样每一轮都只需一次 dead-end 质量求和，不必对每个 dead-end 节点单独发边，既符合冻结公式，又能减少实现复杂度。


与两种对照策略相比：


\begin{itemize}
\item “忽略 dead-end”会导致总概率质量逐渐下降；
\item “删除 dead-end”需要先做递归剪枝，再回填结果，流程更复杂，也会改变原图语义；
\item “质量补偿”既保持原图节点全集，又能维持概率质量守恒，是课程要求下最稳妥的工程方案。
\end{itemize}


\subsection{Top-k 输出与接口对接}


\texttt{dump\_top\_k} 的职责有两层：


\begin{enumerate}
\item 按 \texttt{Score} 降序、\texttt{NodeID} 升序写出 Top-100；
\item 返回 \texttt{top10\_signature}，供 \texttt{benchmark.py} 与 \texttt{sweep.py} 做回归比较。
\end{enumerate}


输出文件每行格式固定为：


\begin{lstlisting}
NodeID Score
\end{lstlisting}


其中 \texttt{Score} 固定保留 10 位小数。由于前面提到的内部重映射，写文件时不能直接把内部数组下标当成 NodeID，而必须先通过 \texttt{\_LAST\_NODE\_IDS} 还原。也正因为如此，实现默认假定同一进程中“最近一次 \texttt{load\_graph} 的图”和“当前要输出的 \texttt{ranks}”是同一份图，这是接口使用时需要注意的一点。


\subsection{本章小结}


本章说明了 PageRank 的 CSR 实现方式。通过“流式读图 + CSR 压缩 + 原地幂迭代 + dead-end 统一补偿”，算法在保持正确性的同时降低了内存占用，并为分块迭代和实验脚本提供了统一接口。

\section{第 4 章 内存优化}


\subsection{问题背景}


PageRank 的核心开销来自两个部分：一是图结构本身的存储，二是每轮幂迭代时对概率向量的更新。若直接采用稠密邻接矩阵，设图有 \texttt{N} 个节点，则需要存储 $N \times N$ 个浮点数。对本作业数据集而言，节点数约为 \texttt{10000}，若使用 \texttt{float32} 存储完整矩阵，光矩阵本体就约为 \texttt{10000} $\times$ \texttt{10000} $\times$ \texttt{4} $\approx$ \texttt{400 MB}，已经远超题目要求的 \texttt{80 MB} 峰值内存限制。因此，本项目必须在数据结构层面先把空间复杂度从 $O(N^2)$ 压到 \texttt{O(M + N)}，再进一步控制迭代过程中的活动工作集。


本章采用两层优化思路。第一层是 CSR（Compressed Sparse Row）稀疏表示，目标是只保存真实存在的边；第二层是按目标节点做真分块，把边先落盘，再用 \texttt{np.memmap} 在每轮按块读取，使 \texttt{r\_new} 的活动更新片段始终限制在当前块内。稀疏矩阵表示和迭代求解是大规模线性系统中的经典思路，相关背景可参考 Saad 与 Davis 的教材\cite{saad2003iterative,davis2006direct}。最终的主提交版本采用 \texttt{CSR + Block Matrix} 组合，其中 CSR 负责压缩图结构，Block Matrix 负责限制迭代阶段的常驻内存。


\begin{figure}[H]
\centering
\safeincludegraphics[width=0.86\textwidth]{implementation_overview.png}
\caption{Dense、CSR 与 CSR+Block 三种实现方式的整体对比。}
\label{fig:implementation-overview}
\end{figure}



\subsection{CSR 的空间优势}


本项目中，A 角色负责输出 CSR 三元组：


\begin{itemize}
\item \texttt{row\_ptr}：长度为 \texttt{N + 1} 的行指针
\item \texttt{col\_idx}：长度为 \texttt{M} 的列索引
\item \texttt{out\_deg}：长度为 \texttt{N} 的出度数组
\end{itemize}


对于源节点 \texttt{u}，其所有出边位于区间：


\begin{lstlisting}
col_idx[row_ptr[u] : row_ptr[u + 1]]
\end{lstlisting}


这样做的好处是只保存真实边，不保存不存在的零元素。空间复杂度从稠密矩阵的 $O(N^2)$ 下降为：


\begin{lstlisting}
O(N + M)
\end{lstlisting}


在本数据集上，\texttt{N = 10000}、\texttt{M = 150000}。如果索引统一使用 \texttt{int32}，则仅就核心结构估算：


\begin{itemize}
\item \texttt{row\_ptr} 约 \texttt{(10001} $\times$ \texttt{4)} $\approx$ \texttt{39 KB}
\item \texttt{col\_idx} 约 \texttt{(150000} $\times$ \texttt{4)} $\approx$ \texttt{586 KB}
\item \texttt{out\_deg} 约 \texttt{(10000} $\times$ \texttt{4)} $\approx$ \texttt{39 KB}
\end{itemize}


图结构本身只需约 \texttt{664 KB} 存储，远低于稠密矩阵方案。同时 CSR 让每轮迭代只需顺着真实边做传播，不需要访问整张 $N \times N$ 矩阵，避免了大量无意义的零乘加。


\subsection{分块策略}


仅有 CSR 还不够。若直接在内存中保留完整的 \texttt{r\_new}、反向邻接关系和辅助临时数组，Python 解释器、NumPy 运行时和中间缓冲区叠加后，RSS 峰值仍可能接近题目上限。因此本项目进一步采用按目标节点分块的策略：


\begin{enumerate}
\item 先按\textbf{目标节点区间}把边分桶。
\item 每个桶落盘为一个独立的 \texttt{.bin} 文件。
\item 每轮迭代时，只 \texttt{memmap} 当前块的边记录。
\item \texttt{r\_new[node\_start:node\_stop]} 作为当前唯一活动更新片段。
\end{enumerate}


按目标节点而不是源节点分块，是因为 PageRank 更新写入的是 \texttt{r\_new[v]}。这样同一个块内的所有边贡献都写回连续的局部区间，避免在多块之间来回写入。


本项目的块文件记录格式为：


\begin{lstlisting}
(src: int32, local_dst: int32)
\end{lstlisting}


其中 \texttt{local\_dst = dst - node\_start}。这样在迭代阶段无需重复做全局下标减法，直接对块内局部数组累加即可。分块构建一次完成，后续每轮复用。


\subsection{memmap 与活动工作集控制}


在分块迭代阶段，完整的 \texttt{r} 与 \texttt{r\_new} 不再作为普通 \texttt{ndarray} 常驻，而是映射到临时文件：


\begin{lstlisting}
r      -> rank_a.bin
r_new  -> rank_b.bin
\end{lstlisting}


这样，操作系统只会把当前访问到的页调入内存。算法本身在代码层面遵循以下约束：


\begin{itemize}
\item \texttt{r\_new} 只通过当前块视图 \texttt{r\_new[node\_start:node\_stop]} 写入
\item 边记录通过 \texttt{np.memmap} 顺序读取
\item 块内贡献使用固定大小缓冲区处理，避免按边数生成大临时数组
\item 每轮结束后只交换 \texttt{r} 与 \texttt{r\_new} 的角色，不复制整向量
\end{itemize}


对应的块迭代公式仍然保持与普通 PageRank 一致：


\begin{lstlisting}
r_new[v] = (1 - beta) / N
         + beta * dangling_mass / N
         + beta * sum(r[u] / out_deg[u] for u -> v)
\end{lstlisting}


区别在于”对每个 \texttt{v} 的更新”变成了”对每个目标块的更新”，分块不改变数值结果，只改变访问顺序和内存布局。


\begin{lstlisting}
伪代码：
for iter in range(max_iter):
    dangling_mass = sum(r[u] for u in dead_ends)
    base = (1 - beta) / N + beta * dangling_mass / N
    for each block:
        block_view.fill(base)
        for each edge (src, local_dst) in block_file:
            block_view[local_dst] += beta * r[src] / out_deg[src]
    check L1 convergence
    swap(r, r_new)
\end{lstlisting}


\subsection{float32 的选择}


所有概率向量默认采用 \texttt{float32}。PageRank 本身是迭代近似问题，题目规定的停止阈值为 \texttt{1e-8}，不需要双精度的极限精度；\texttt{float32} 可以把向量和缓冲区的空间消耗减半，而本项目的评价重点在于时间与 RSS 峰值。


实验 E7 显示，\texttt{float32} 与 \texttt{float64} 在 Top-10 签名上保持一致，说明排序层面没有发生可见偏差。3 次重复平均结果如下：


\begin{itemize}
\item \texttt{float32}：峰值 RSS 约 \texttt{35.92 MB}，平均耗时约 \texttt{0.403 s}，迭代 \texttt{15} 轮
\item \texttt{float64}：峰值 RSS 约 \texttt{36.27 MB}，平均耗时约 \texttt{0.366 s}，迭代 \texttt{12} 轮
\end{itemize}


\texttt{float64} 耗时略低，主要是因为收敛轮数从 15 降到了 12，而不是单次运算更快。不过题目把内存控制作为核心约束，主提交版本仍固定为 \texttt{float32}，在更大数据规模上能留出更多安全余量。


\subsection{实测结果与分析}


在 \texttt{Data.txt} 数据集、\texttt{beta = 0.85}、\texttt{eps = 1e-8} 条件下，实验 E3 表明块数 \texttt{K} 会在“块内数组大小”和“块切换开销”之间形成权衡：


\begin{itemize}
\item \texttt{K = 4}：平均峰值 RSS 约 \texttt{35.16 MB}，平均时间约 \texttt{0.376 s}
\item \texttt{K = 8}：平均峰值 RSS 约 \texttt{35.99 MB}，平均时间约 \texttt{0.423 s}
\item \texttt{K = 16}：平均峰值 RSS 约 \texttt{34.97 MB}，平均时间约 \texttt{0.457 s}
\end{itemize}


在当前数据规模下，增大 \texttt{K} 并不会显著降低 RSS，反而增加了块切换和文件访问开销。\texttt{N = 10000}、\texttt{M = 150000} 本身并不大，Python 解释器、NumPy 运行时和子进程开销已经占了相当比例，过细分块的收益有限。


E4 对比纯 CSR 与 \texttt{CSR + Block} 的结果如下：


\begin{itemize}
\item \texttt{csr}：平均峰值 RSS 约 \texttt{33.42 MB}，平均时间约 \texttt{0.492 s}
\item \texttt{csr\_block}：平均峰值 RSS 约 \texttt{35.89 MB}，平均时间约 \texttt{0.418 s}
\end{itemize}


分块版本在时间上略有优势，但峰值 RSS 稍高。在当前数据规模上，分块更多体现为可扩展的内存控制手段；若数据集扩大，分块的收益会更明显，因为活动工作集不再随全部边数线性增长。


\subsection{本章结论}


内存优化分为几层：CSR 把空间复杂度从 $O(N^2)$ 压到 $O(N + M)$；分块 + \texttt{memmap} 把迭代时的活动工作集限制在当前块；\texttt{float32} 进一步降低向量和缓冲区的内存占用。最终主提交版本 \texttt{csr\_block + float32 + K=8} 的平均峰值 RSS 约 \texttt{36 MB}，低于作业要求的 \texttt{80 MB} 上限。

\section{第 5 章 实验结果与分析}


\subsection{实验设计与数据来源}


本章汇总数据集统计、分块参数、阻尼系数、收敛阈值和浮点精度实验。所有图表均由脚本从 CSV 生成，不手工填写图中数据。数据来源如下：


\begin{itemize}
\item 数据集分析：\texttt{experiments/dataset\_stats.json} 与 \texttt{experiments/degree\_distribution.csv}
\item E3、E4、E7：\texttt{experiments/E3.csv}、\texttt{experiments/E4.csv}、\texttt{experiments/E7.csv}
\item E5、E6：\texttt{sweep.py} 通过 \texttt{main.py} CLI 生成的 \texttt{experiments/sweep\_results.csv}
\item 图表输出：\texttt{report/fig/*.png}
\end{itemize}


\subsection{数据集分析结果}


数据集共有 \texttt{10000} 个节点、\texttt{150000} 条边，平均入度与平均出度均为 \texttt{15.00}。最大出度为 \texttt{36}，最大入度为 \texttt{32}，说明图整体较稀疏，且不存在极端高连接度的超级节点。若采用稠密矩阵，空间复杂度为 $O(N^2)$，对一万节点图会形成数百 MB 的基础内存占用；而 CSR 将图结构压缩为 \texttt{row\_ptr / col\_idx / out\_deg} 三类数组，空间复杂度降为 \texttt{O(N + M)}。


数据集中 dead-end 节点数为 \texttt{1500}，占比 \texttt{15.00\%}。这是算法实现中最需要关注的结构特征之一。如果不补偿这些节点持有的概率质量，PageRank 向量的总和会逐轮下降，最终得到的排名不再对应合法概率分布。主程序和分块版本均采用相同的质量补偿公式。


弱连通分量共有 \texttt{501} 个，最大弱连通分量包含 \texttt{9500} 个节点，占比 \texttt{95.00\%}，另有 \texttt{500} 个孤立节点。这说明图主体结构集中，Top 排名大概率来自最大弱连通分量；孤立节点和小分量则主要依赖 teleport 与 dead-end 补偿获得基础分数。


\subsection{E3：块数 K 对内存与时间的影响}


对 \texttt{csr\_block} 模式在 \texttt{K = 4 / 8 / 16} 下各运行 3 次，平均结果如下：


\begin{center}
\small
\begin{tabularx}{\textwidth}{>{\raggedright\arraybackslash}X >{\raggedright\arraybackslash}X >{\raggedright\arraybackslash}X >{\raggedright\arraybackslash}X}
\toprule
K & 平均峰值 RSS / MB & 平均时间 / s & 迭代轮数 \\
\midrule
4 & 35.16 & 0.376 & 15 \\
8 & 35.99 & 0.423 & 15 \\
16 & 34.97 & 0.457 & 15 \\
\bottomrule
\end{tabularx}
\end{center}


图~\ref{fig:k-tradeoff} 展示了块数与内存、时间之间的权衡。增大 \texttt{K} 并没有显著降低内存峰值，\texttt{K=16} 的 RSS 只比 \texttt{K=8} 略低，但耗时更高。原因在于本数据集只有十五万条边，Python 解释器、NumPy 运行时和进程采样开销已经占据相当比例；当图本身不大时，过细分块会增加文件切换和块循环成本。


\begin{figure}[H]
\centering
\safeincludegraphics[width=0.86\textwidth]{k_tradeoff.png}
\caption{块数 $K$ 对峰值内存和运行时间的影响。}
\label{fig:k-tradeoff}
\end{figure}



因此，\texttt{K=8} 可作为默认设置：运行时间与内存均处于安全范围内，也为更大图保留了分块扩展能力。


\subsection{E4：CSR 与 CSR+Block 对比}


实验 E4 对比了纯 \texttt{csr} 与 \texttt{csr\_block}，平均结果如下：


\begin{center}
\small
\begin{tabularx}{\textwidth}{>{\raggedright\arraybackslash}X >{\raggedright\arraybackslash}X >{\raggedright\arraybackslash}X >{\raggedright\arraybackslash}X}
\toprule
模式 & 平均峰值 RSS / MB & 平均时间 / s & 迭代轮数 \\
\midrule
\texttt{csr} & 33.42 & 0.492 & 15 \\
\texttt{csr\_block} & 35.89 & 0.418 & 15 \\
\bottomrule
\end{tabularx}
\end{center}


就结果看，\texttt{csr\_block} 的时间略优于纯 CSR，但 RSS 略高。这说明分块收益会被 Python 基线开销和 \texttt{memmap} 管理开销抵消一部分；其主要价值在于限制活动更新区，使更大规模图上的内存上界更可控。


\subsection{E5：阻尼系数 beta 对 Top-10 排名的影响}


使用 \texttt{sweep.py} 扫描 $\beta \in \{0.70, 0.80, 0.85, 0.90, 0.95\}$，固定 \texttt{eps=1e-8}、\texttt{mode=csr\_block}、\texttt{K=8}、\texttt{dtype=float32}。以 \texttt{beta=0.85} 的 Top-10 为基准，计算 Jaccard 相似度和交集上的 Kendall tau，结果如下：


\begin{center}
\small
\begin{tabularx}{\textwidth}{
  >{\raggedright\arraybackslash}p{1.0cm}
  >{\raggedright\arraybackslash}p{1.4cm}
  >{\raggedright\arraybackslash}p{2.0cm}
  >{\raggedright\arraybackslash}p{1.8cm}
  >{\raggedright\arraybackslash}X
}
\toprule
beta & 迭代轮数 & Top-10 Jaccard & Kendall tau & Top-10 签名 \\
\midrule
0.70 & 13 & 0.818 & 0.833 &
\makecell[l]{\ttfamily 75,8686,5104,725,3257,\\9678,468,7730,7175,6165} \\

0.80 & 15 & 1.000 & 0.956 &
\makecell[l]{\ttfamily 75,8686,5104,9678,725,\\3257,468,7730,7175,5526} \\

0.85 & 15 & 1.000 & 1.000 &
\makecell[l]{\ttfamily 75,8686,9678,5104,725,\\3257,468,7730,7175,5526} \\

0.90 & 16 & 1.000 & 0.956 &
\makecell[l]{\ttfamily 75,8686,9678,5104,725,\\3257,468,7175,7730,5526} \\

0.95 & 16 & 0.818 & 0.944 &
\makecell[l]{\ttfamily 75,8686,9678,5104,725,\\468,3257,7175,5526,8880} \\
\bottomrule
\end{tabularx}
\end{center}


图~\ref{fig:beta-sensitivity} 显示，\texttt{beta} 在 \texttt{0.80} 到 \texttt{0.90} 区间时 Top-10 集合非常稳定，仅发生局部顺序变化；当 \texttt{beta=0.70} 或 \texttt{0.95} 时，Top-10 集合出现 1 个节点替换。较低 beta 会增强随机跳转项并削弱链接结构，较高 beta 会强化沿边传播并放大局部结构差异。


\begin{figure}[H]
\centering
\safeincludegraphics[width=0.86\textwidth]{beta_sensitivity.png}
\caption{阻尼系数 $\beta$ 对 Top-10 稳定性和迭代轮数的影响。}
\label{fig:beta-sensitivity}
\end{figure}



\subsection{E6：收敛阈值 eps 对结果与迭代轮数的影响}


使用 \texttt{sweep.py} 扫描 $\varepsilon \in \{\texttt{1e-6}, \texttt{1e-8}, \texttt{1e-10}\}$，固定 \texttt{beta=0.85}、\texttt{mode=csr\_block}、\texttt{K=8}、\texttt{dtype=float32}。结果如下：


\begin{center}
\small
\begin{tabularx}{\textwidth}{>{\raggedright\arraybackslash}X >{\raggedright\arraybackslash}X >{\raggedright\arraybackslash}X >{\raggedright\arraybackslash}X}
\toprule
eps & 迭代轮数 & 最终 L1 delta & Top-10 Jaccard \\
\midrule
\texttt{1e-6} & 9 & \texttt{3.503e-7} & 1.000 \\
\texttt{1e-8} & 15 & \texttt{7.170e-9} & 1.000 \\
\texttt{1e-10} & 27 & \texttt{8.004e-11} & 1.000 \\
\bottomrule
\end{tabularx}
\end{center}


图~\ref{fig:epsilon-sensitivity} 展示了阈值变严格后迭代轮数增加的趋势。\texttt{eps=1e-6} 到 \texttt{1e-8} 时，轮数从 9 增至 15；进一步到 \texttt{1e-10} 时，轮数增至 27。三组实验的 Top-10 集合和顺序均与 \texttt{eps=1e-8} 保持一致，说明在本数据集上，Top-10 排名对收敛阈值不敏感。


\begin{figure}[H]
\centering
\safeincludegraphics[width=0.86\textwidth]{epsilon_sensitivity.png}
\caption{收敛阈值 $\varepsilon$ 对迭代轮数和最终残差的影响。}
\label{fig:epsilon-sensitivity}
\end{figure}



本实验记录最终 L1 residual 与迭代轮数，用于比较不同收敛阈值下的代价与排名稳定性。


\subsection{E7：dtype 对比}


实验 E7 对比了 \texttt{float32} 与 \texttt{float64}：


\begin{center}
\small
\begin{tabularx}{\textwidth}{>{\raggedright\arraybackslash}X >{\raggedright\arraybackslash}X >{\raggedright\arraybackslash}X >{\raggedright\arraybackslash}X}
\toprule
dtype & 平均峰值 RSS / MB & 平均时间 / s & 迭代轮数 \\
\midrule
\texttt{float32} & 35.92 & 0.403 & 15 \\
\texttt{float64} & 36.27 & 0.366 & 12 \\
\bottomrule
\end{tabularx}
\end{center}


\begin{figure}[H]
\centering
\safeincludegraphics[width=0.86\textwidth]{dtype_comparison.png}
\caption{\texttt{float32} 与 \texttt{float64} 的内存和时间对比。}
\label{fig:dtype-comparison}
\end{figure}


\texttt{float64} 因收敛轮数较少，平均时间略低；但它的内存峰值高于 \texttt{float32}。考虑到内存约束，主提交版本固定为 \texttt{float32}。从 Top-10 签名看，两种 dtype 未造成可见排名差异。


\subsection{实验结论}


综合来看，数据集稀疏性明显，CSR 是满足内存限制的必要选择。dead-end 占比 \texttt{15.00\%}，必须统一做质量补偿。\texttt{beta} 在 \texttt{0.80} 到 \texttt{0.90} 区间时 Top-10 排名稳定，极端值会引入小幅变化。收敛阈值越严格迭代轮数越多，但 Top-10 排名在本数据集上对 \texttt{eps} 不敏感。


\section{第 6 章 工程实现与部署}

\subsection{工程目标}

工程部分的目标是保证 PageRank 实现可以稳定运行、方便复现实验，并能按统一格式输出结果。主入口为 \texttt{main.py}，负责参数解析、算法调度、结果写入和运行摘要输出；分块迭代由 \texttt{blocks.py} 提供，实验统计与绘图脚本统一放在 \texttt{scripts/} 目录中，避免根目录混杂临时工具。

\subsection{项目结构}

项目文件按功能组织为三类：

\begin{itemize}
\item 核心实现：\texttt{main.py}、\texttt{blocks.py}、\texttt{mock\_graph.py}；
\item 实验工具：\texttt{scripts/benchmark.py}、\texttt{scripts/sweep.py}、\texttt{scripts/plot.py}、\texttt{scripts/run\_e8.py}；
\item 辅助工具：\texttt{scripts/memoryuse-python.py} 保留老师提供的内存采样脚本，用于人工复核；
\item 配置与说明：\texttt{requirements.txt}、\texttt{docs/packaging/compile-parameter.txt}、\texttt{README.md}。
\end{itemize}

这种结构使算法实现、分块优化和实验绘图相互解耦。实验脚本只通过命令行调用主程序，不直接依赖内部函数，从而降低接口变动带来的影响。

\subsection{命令行接口}

主程序采用统一命令行参数：

\begin{lstlisting}
python main.py --data Data.txt --beta 0.85 --eps 1e-8 --out Res.txt --mode csr_block --K 8 --dtype float32
\end{lstlisting}

其中 \texttt{--mode} 用于选择 \texttt{dense / csr / block / csr\_block}，\texttt{--K} 控制分块数，\texttt{--dtype} 控制排名向量精度。主提交模式为 \texttt{csr\_block + K=8 + float32}。

程序运行结束后，除生成 \texttt{Res.txt} 外，还会在标准输出末尾给出 JSON 摘要：

\begin{lstlisting}
{
  "peak_rss_mb": 35.6,
  "wall_sec": 0.38,
  "iters": 15,
  "mode": "csr_block",
  "K": 8,
  "dtype": "float32",
  "top10_signature": "75,8686,..."
}
\end{lstlisting}

该摘要便于 \texttt{scripts/benchmark.py} 自动收集运行时间、内存峰值、迭代轮数和 Top-10 签名。

\subsection{正确性与复现}

正确性验证包括两类测试。第一类是小图测试，用于检查 dead-end 补偿、spider-trap 抑制和收敛判据是否符合公式。第二类是随机图对拍，用于比较普通 CSR 迭代和分块迭代的结果是否一致。实验结果统一写入 CSV，再由绘图脚本生成图表，保证报告中的数据可追溯。

开发与复现主要在 Windows 的 \texttt{micromamba activate test} 环境中完成；同时使用 WSL 调用 Windows 版 micromamba 复跑源码测试，确认路径和命令在挂载目录下可用。由于 PowerShell 会把未加引号的逗号参数拆分，批量实验命令中的模式列表写为 \texttt{--modes 'csr,csr\_block'}。

\subsection{打包与运行}

最终可使用 PyInstaller 将 \texttt{main.py} 打包为单文件可执行程序，打包参数保存在 \texttt{docs/packaging/compile-parameter.txt}。打包前先运行源码版本，确认 \texttt{Res.txt}、stdout JSON 和 Top-10 签名均正常；打包后再使用相同参数复测，比较源码版与可执行版的输出是否一致。该流程主要验证路径解析、依赖收集和结果格式，不改变算法逻辑。

为便于换机检查，项目额外准备了本地可移动验证包 \texttt{portable\_check/}。该目录包含 \texttt{main.exe}、\texttt{Data.txt}、\texttt{run\_check.ps1} 和简短说明；PowerShell 脚本会运行主提交模式，生成 \texttt{Res\_check.txt}，检查结果文件行数和 Top-10 签名。由于其中包含数据集，该目录仅作为本地验证工具，不默认加入最终提交压缩包。

\subsection{本章小结}

工程实现围绕统一入口、固定输出格式和可复现实验展开。\texttt{main.py} 提供稳定 CLI，\texttt{scripts/benchmark.py} 与 CSV 记录实验数据，PyInstaller 用于最终交付，\texttt{portable\_check/} 用于换机复核。这样既保证 PageRank 结果可复现，也便于比较 CSR、分块数和浮点精度等不同设置的影响。


\section{第 7 章 总结与反思}

\subsection{工作总结}

本项目完成了 PageRank 在 \texttt{Data.txt} 有向图上的计算，并输出 Top-100 节点排名。实现上，算法使用 CSR 表示稀疏图，用 dead-end 质量补偿保证概率守恒，用 teleport 项抑制 spider-trap，并通过幂迭代求解稳定分布。在内存优化方面，项目进一步引入按目标节点分块和 \texttt{memmap}，将迭代阶段的活动工作集限制在当前块内。

实验表明，\texttt{csr\_block + K=8 + float32} 能在该数据集上稳定收敛，Top-10 结果在常用阻尼系数和收敛阈值范围内保持较高稳定性，峰值内存也低于题目限制。

\subsection{团队分工}

小组在分工时优先保证工作量均衡，A/B/C 角色通过抽签方式确定。角色字母仅用于区分算法核心、工程优化和实验分析三类职责，不表示贡献权重高低；最终报告、代码和实验结果均由小组成员共同复核。

\begin{center}
\small
\begin{tabularx}{\textwidth}{>{\raggedright\arraybackslash}p{2.7cm} >{\raggedright\arraybackslash}p{2.8cm} >{\raggedright\arraybackslash}X >{\raggedright\arraybackslash}X}
\toprule
成员 & 实验与章节 & 代码/脚本 & 主要职责 \\
\midrule
匡航逸（A） & E1、E2、E8；第 2、3 章 & \texttt{main.py} 核心接口与 CSR 迭代 & PageRank 核心算法与正确性验证 \\
柯云超（B） & E3、E4、E7；第 4、6 章 & \texttt{blocks.py} 与 \texttt{scripts/benchmark.py} & 分块优化、内存测试与工程集成 \\
蒋林瀞（C） & E5、E6、数据集分析；第 1、5、7 章 & \texttt{scripts/} 数据分析、扫描与绘图脚本 & 数据分析、参数实验与报告整理 \\
\bottomrule
\end{tabularx}
\end{center}

\subsection{经验与反思}

PageRank 的关键不只是入边统计，而是概率质量在全图上的稳定分布。dead-end 处理不当会造成质量泄漏，spider-trap 缺少 teleport 会造成局部吸收，两者必须在公式和实现中同时考虑。

内存优化需要区分理论数组空间和实际 RSS。CSR 能降低图结构存储，但 Python 解释器、NumPy 运行时和临时缓冲区也会影响峰值内存，报告中的内存结论应以 benchmark 实测为准。

分块优化在小图上不一定带来最低 RSS，但提供了更清晰的内存上界控制方式。图规模扩大时，按目标节点分块可以避免一次性维护过大的更新工作集。

本项目完成了从 PageRank 理论公式到稀疏存储、分块迭代、实验验证和工程交付的完整流程。实验结果表明 CSR 与 Block Matrix 的结合能够在保持结果一致性的同时满足内存约束。

\subsection{AI 工具使用说明}

项目整理和报告定稿过程中使用了 OpenAI Codex 辅助检查目录结构和整理文档表述。所有 PageRank 算法实现、实验数据和最终结论均由小组成员人工复核，代码未调用任何现成 PageRank API，报告中的实验数字来自项目脚本和运行日志。


\renewcommand{\refname}{参考文献}
\nocite{*}
\bibliographystyle{plain}
\bibliography{reference}
\end{document}
